%%%%%%%%%%%%%%%%%%%%%%%%
%
% $Autor: Hemanth Jadiswami Prabhakaran $
% $Datum: 2026-09-25 $
% $Pfad: GitHub/MBIDA_Master_Thesis/Manual/Chapters/en/04UserInterface.tex $
% $Version: 1 $
%
% $Project: MBIDA Thesis - Manual $
%
%%%%%%%%%%%%%%%%%%%%%%%%


\chapter{User Interface}
\label{chap:ui}

The dashboard consists of a single page. There are no menus to open, no tabs and no pop-up windows: every control is visible at all times, and every change takes effect immediately. This chapter first presents a sketch of the complete screen and then describes its two parts, the sidebar and the chart.

\section{Sketch of the User Interface}
\label{sec:uiSketch}

Figure~\ref{fig:uiWireframe} is a schematic sketch of the screen. The numbered elements are listed in Table~\ref{tab:uiCallouts}; the same numbers are used throughout the manual.

\begin{figure}[H]
  \centering
  \resizebox{\linewidth}{!}{% GUI sketch of the dashboard with numbered callouts.
% The numbers are explained in Table "tab:uiCallouts" (Chapter 4).
\begin{tikzpicture}[
    font=\sffamily,
    lbl/.style={anchor=west,font=\sffamily\scriptsize,text=AppText},
    hdr/.style={anchor=west,font=\sffamily\small\bfseries,text=AppText},
    sel/.style={draw=ManGrey!60,fill=white,rounded corners=2pt,minimum width=3.7cm,
                minimum height=0.48cm,anchor=west,font=\ttfamily\scriptsize,text=AppText,
                inner sep=3pt,text width=3.3cm},
  ]

  % ------------------------------------------------ page + sidebar
  \draw[ManGrey!70,fill=white,rounded corners=3pt] (0,0) rectangle (16,11);
  \fill[AppSidebar,rounded corners=3pt] (0,0) rectangle (4.3,11);
  \fill[AppSidebar] (3.0,0) rectangle (4.3,11);

  % collapse chevron + page menu
  \node[font=\sffamily\small,text=ManGrey] at (4.0,10.6) {\guillemotleft};
  \node[font=\sffamily\bfseries,text=ManGrey,rotate=90] at (15.6,10.6) {\textbf{$\cdots$}};

  % ------------------------------------------------ sidebar: hierarchy
  \node[hdr] at (0.25,10.05) {Hierarchy};
  \foreach \name/\val/\y in {State/CA/9.35, Store/CA\_1/8.3, Category/(All)/7.25,
                             Department/(All)/6.2, Item/(All)/5.15}{
     \node[lbl] at (0.25,\y) {\name};
     \node[sel] (s\name) at (0.3,\y-0.43) {\val};
     \node[font=\tiny,text=ManGrey] at (3.8,\y-0.43) {$\blacktriangledown$};
  }

  % ------------------------------------------------ sidebar: forecast
  \node[hdr] at (0.25,4.05) {Forecast};
  \node[lbl] at (0.25,3.5) {Base model};
  \foreach \x/\m/\on in {0.45/histavg/0, 1.75/ets/1, 2.75/theta/0}{
     \ifnum\on=1
       \draw[AppRed,fill=AppRed] (\x,3.05) circle (0.1);
       \fill[white] (\x,3.05) circle (0.04);
     \else
       \draw[ManGrey] (\x,3.05) circle (0.1);
     \fi
     \node[lbl] at (\x+0.1,3.05) {\m};
  }
  \node[lbl] at (0.25,2.5) {Reconciliation method(s)};
  \node[sel,minimum height=0.62cm] at (0.3,1.95) {};
  \node[anchor=west,fill=AppRed,text=white,rounded corners=2pt,font=\sffamily\tiny,
        inner sep=2.5pt] at (0.4,1.95) {MinTrace (wls\_struct) \,$\times$};

  % ------------------------------------------------ main area: header
  \node[anchor=west,font=\sffamily\large\bfseries,text=AppText] at (4.8,10.1)
       {M5 Hierarchical Forecast Explorer};
  \node[anchor=west,font=\sffamily\tiny,text=ManGrey,text width=10.4cm] at (4.8,9.45)
       {Drill down Total $\rightarrow$ State $\rightarrow$ Store $\rightarrow$ Category
        $\rightarrow$ Department $\rightarrow$ Item. Leave a level on ``(All)'' to view
        that aggregate node instead of a single item. Click a dropdown and type to search.};
  \node[anchor=west,font=\sffamily\scriptsize,text=AppText] at (4.8,8.8)
       {Selected node: \textbf{Total/CA/CA\_1} \ (level = 3)};

  % ------------------------------------------------ chart
  \pgfmathsetmacro{\xa}{5.3}\pgfmathsetmacro{\xb}{15.5}
  \pgfmathsetmacro{\ya}{1.0}\pgfmathsetmacro{\yb}{7.55}
  \pgfmathsetmacro{\sx}{(\xb-\xa)/66}
  \pgfmathsetmacro{\xtest}{\xa+47*\sx}
  \pgfmathsetmacro{\xlive}{\xa+59*\sx}

  \fill[AppTrain] (\xa,\ya) rectangle (\xtest,\yb);
  \fill[AppTest]  (\xtest,\ya) rectangle (\xlive,\yb);
  \fill[AppLive]  (\xlive,\ya) rectangle (\xb,\yb);
  \foreach \x/\t in {\xa/Train,\xtest/Test,\xlive/Live}
     \node[anchor=north west,font=\sffamily\tiny,text=ManGrey] at (\x,\yb) {\t};

  % grid + axes
  \foreach \k in {1,...,4}
     \draw[ManGrey!25] (\xa,{\ya+\k*1.3}) -- (\xb,{\ya+\k*1.3});
  \draw[ManGrey!60] (\xa,\ya) -- (\xb,\ya);
  \foreach \m/\t in {11/2012,23/2013,35/2014,47/2015,59/2016}
     \node[font=\sffamily\tiny,text=ManGrey,below] at ({\xa+\m*\sx},\ya) {\t};
  \node[font=\sffamily\tiny,text=AppText] at ({(\xa+\xb)/2},0.35) {Date};
  \node[font=\sffamily\tiny,text=AppText,rotate=90] at (5.0,{(\ya+\yb)/2}) {Sales};

  % traces (illustrative shapes)
  \draw[AppActual,thick] plot[domain=0:58,samples=59,mark=*,mark size=0.5pt]
       ({\xa+\x*\sx},{3.0+0.028*\x+0.55*sin(30*\x)+0.9*exp(-((mod(\x,12)-10)^2))});
  \draw[AppBase,thick,dotted] plot[domain=47:65,samples=19]
       ({\xa+\x*\sx},{2.75+0.028*\x+0.45*sin(30*\x)+0.6*exp(-((mod(\x,12)-10)^2))});
  \draw[AppRecA,thick] plot[domain=47:65,samples=19]
       ({\xa+\x*\sx},{2.95+0.028*\x+0.5*sin(30*\x)+0.8*exp(-((mod(\x,12)-10)^2))});

  % legend (horizontal, above the plot, right aligned)
  \draw[AppActual,thick] (8.3,8.3) -- (8.7,8.3);  \fill (8.5,8.3) circle (0.04);
  \node[lbl,font=\sffamily\tiny] at (8.72,8.3) {Actual};
  \draw[AppBase,thick,dotted] (9.75,8.3) -- (10.15,8.3);
  \node[lbl,font=\sffamily\tiny] at (10.17,8.3) {Base: ets};
  \draw[AppRecA,thick] (11.45,8.3) -- (11.85,8.3);
  \node[lbl,font=\sffamily\tiny] at (11.87,8.3) {Reconciled: MinTrace (wls\_struct)};

  % Plotly toolbar (appears on hover)
  \foreach \k in {0,...,6}
     \draw[ManGrey!70,fill=white,rounded corners=1pt] ({14.0+\k*0.22},7.8) rectangle ++(0.17,0.17);

  % ------------------------------------------------ callouts
  \node[man/callout] at (3.55,10.6) {1};
  \draw[ManRed!70,decorate,decoration={brace,amplitude=4pt,mirror}] (-0.15,9.55) -- (-0.15,4.45);
  \node[man/callout] at (-0.55,7.0) {2};
  \node[man/callout] at (-0.35,3.05) {3};
  \node[man/callout] at (-0.35,1.95) {4};
  \node[man/callout] at (4.55,9.8) {5};
  \node[man/callout] at (4.55,8.8) {6};
  \node[man/callout] at (8.0,8.3) {7};
  \node[man/callout] at (\xtest,6.9) {8};
  \node[man/callout] at (7.5,5.6) {9};
  \node[man/callout] at (13.7,7.88) {10};
  \node[man/callout] at (15.05,10.6) {11};
  \node[man/callout] at ({(\xa+\xb)/2+0.6},0.35) {12};

\end{tikzpicture}
}
  \caption{Sketch of the user interface with numbered elements}
  \label{fig:uiWireframe}
\end{figure}

\begin{table}[H]
\centering
\caption{Elements of the user interface (numbers as in Figure~\ref{fig:uiWireframe})}
\label{tab:uiCallouts}
\small
\begin{tabularx}{\linewidth}{@{}clX@{}}
\toprule
\textbf{No.} & \textbf{Element} & \textbf{Purpose} \\
\midrule
\Callout{1}  & Sidebar button        & Hides or shows the sidebar, giving the chart the full width \\
\Callout{2}  & Hierarchy dropdowns   & Select the node: \UI{State}, \UI{Store}, \UI{Category}, \UI{Department}, \UI{Item} \\
\Callout{3}  & \UI{Base model}       & Choose \VAL{histavg}, \VAL{ets} or \VAL{theta} for the dotted line \\
\Callout{4}  & \UI{Reconciliation method(s)} & Choose one, several or no reconciled lines \\
\Callout{5}  & Title and instructions & Name of the dashboard and a two-line usage hint \\
\Callout{6}  & Node line              & The node currently shown and its level (1--6) \\
\Callout{7}  & Legend                 & Names of all lines; click to hide or show a line \\
\Callout{8}  & Time regions           & Shaded \UI{Train}, \UI{Test} and \UI{Live} periods \\
\Callout{9}  & Plot area              & Actual sales, base forecast and reconciled forecast(s) \\
\Callout{10} & Chart toolbar          & Save as image, zoom, pan, reset (appears on mouse-over) \\
\Callout{11} & Page menu              & Streamlit's menu: theme, rerun, print \\
\Callout{12} & Axes                   & Month (horizontal) and units sold (vertical) \\
\bottomrule
\end{tabularx}
\end{table}

Figure~\ref{fig:uiReal} shows the same screen as it appears in the browser.

\begin{figure}[H]
  \centering
  \Screenshot{Images/04UserInterface/FullScreen.png}
  \caption{The real user interface for store \VAL{CA\_1} with base model \VAL{ets}}
  \label{fig:uiReal}
\end{figure}

\section{Layout Principles}
\label{sec:uiPrinciples}

\begin{itemize}
  \item \textbf{Controls left, result right.} Everything the user can change is in the sidebar; everything the dashboard shows is in the main area.
  \item \textbf{Top to bottom equals coarse to fine.} The dropdowns follow the hierarchy from State down to Item, so the sidebar is read in the same order as the hierarchy.
  \item \textbf{One chart, one node.} The chart always shows exactly one node. Comparing nodes is done by switching between them.
  \item \textbf{Always consistent.} A dropdown only offers values that exist below the choices made above it, so an impossible combination cannot be selected.
\end{itemize}

\section{The Sidebar}
\label{sec:uiSidebar}

\subsection{Hierarchy Section}
\label{sec:uiHierarchy}

The upper part of the sidebar \Callout{2} holds one dropdown per hierarchy level (Figure~\ref{fig:uiSidebar}, left).

\begin{figure}[H]
  \centering
  \Screenshot[0.42\linewidth]{Images/04UserInterface/SidebarHierarchy.png}
  \hfill
  \Screenshot[0.42\linewidth]{Images/04UserInterface/SidebarForecast.png}
  \caption{Sidebar: hierarchy section (left) and forecast section (right)}
  \label{fig:uiSidebar}
\end{figure}

The five dropdowns are linked:

\begin{table}[H]
\centering
\caption{Hierarchy dropdowns}
\label{tab:uiDropdowns}
\small
\begin{tabularx}{\linewidth}{@{}llX@{}}
\toprule
\textbf{Dropdown} & \textbf{Values} & \textbf{Available when} \\
\midrule
\UI{State}      & \VAL{(All)}, \VAL{CA}, \VAL{TX}, \VAL{WI} & always \\
\UI{Store}      & \VAL{(All)} + stores of the chosen state, e.g.\ \VAL{CA\_1} to \VAL{CA\_4} & a state is chosen \\
\UI{Category}   & \VAL{(All)}, \VAL{FOODS}, \VAL{HOBBIES}, \VAL{HOUSEHOLD} & a store is chosen \\
\UI{Department} & \VAL{(All)} + departments of the category, e.g.\ \VAL{FOODS\_1} to \VAL{FOODS\_3} & a category is chosen \\
\UI{Item}       & \VAL{(All)} + items of the department in that store & a department is chosen \\
\bottomrule
\end{tabularx}
\end{table}

As long as a dropdown is on \VAL{(All)}, all dropdowns below it offer only \VAL{(All)}. Each dropdown can be searched by typing (Section~\ref{sec:fnSearch}).

\subsection{Forecast Section}
\label{sec:uiForecast}

The lower part of the sidebar holds the two forecast controls \Callout{3} and \Callout{4} (Figure~\ref{fig:uiSidebar}, right).

\begin{description}[style=nextline,leftmargin=0.5cm,font=\normalfont\bfseries]
  \item[\UI{Base model} \Callout{3}] Three round buttons in one row. Exactly one is active; the default is \VAL{histavg}.
  \item[\UI{Reconciliation method(s)} \Callout{4}] A multi-selection box. Selected methods appear as red tags with a $\times$ to remove them. Clicking into the box lists the remaining methods. The default is \VAL{MinTrace (wls\_struct)}. The box may also be left empty.
\end{description}

\subsection{Collapsing the Sidebar}

On small screens the sidebar takes a lot of space. The double arrow \Callout{1} at its top right hides it; the chart then uses the full window width. A small arrow at the top left of the page brings the sidebar back. On phones the sidebar is hidden by default.

\section{The Main Area}
\label{sec:uiMain}

\subsection{Title, Instructions and Node Line}

Below the title \Callout{5}, a short grey text summarises how to use the dropdowns. The line directly above the chart, the \emph{node line} \Callout{6}, always states which node is shown:

\begin{lstlisting}[style=plain]
Selected node: Total/CA/CA_1  (level = 3)
\end{lstlisting}

The node name lists the path through the hierarchy, separated by slashes; the level counts from 1 (grand total) to 6 (single item). If a selected reconciliation method cannot be found in the result files, a warning is appended to this line (Section~\ref{sec:tsNoColumn}).

\subsection{The Forecast Chart}
\label{sec:uiChart}

Figure~\ref{fig:chartAnatomy} explains how to read the plot area \Callout{9} with its time regions \Callout{8} and axes \Callout{12}. The data in the sketch is illustrative; the structure is exactly that of the real chart.

\begin{figure}[H]
  \centering
  % Anatomy of the forecast chart - illustrative data, not real results.
% x = month index, 0 = Feb 2011, 47 = Jan 2015 (test starts), 59 = Jan 2016 (live starts).
\begin{tikzpicture}
  \begin{axis}[
      width=\linewidth, height=7.2cm,
      xmin=0, xmax=65, ymin=60, ymax=175,
      xtick={11,23,35,47,59}, xticklabels={2012,2013,2014,2015,2016},
      ytick=\empty,
      xlabel={Date}, ylabel={Sales},
      axis lines*=left,
      tick label style={font=\sffamily\scriptsize},
      label style={font=\sffamily\small},
      clip=false,
      declare function={
        season(\x)=12*sin(30*\x)+22*exp(-((mod(\x,12)-10)^2));
        act(\x)=95+0.55*\x+season(\x);
        base(\x)=90+0.55*\x+0.7*season(\x);
        rec(\x)=94+0.55*\x+0.9*season(\x);
      },
    ]
    % shaded phases
    \fill[AppTrain] (axis cs:0,60)  rectangle (axis cs:47,175);
    \fill[AppTest]  (axis cs:47,60) rectangle (axis cs:59,175);
    \fill[AppLive]  (axis cs:59,60) rectangle (axis cs:65,175);
    \node[anchor=north west,font=\sffamily\scriptsize,text=ManGrey] at (axis cs:0,175)  {Train};
    \node[anchor=north west,font=\sffamily\scriptsize,text=ManGrey] at (axis cs:47,175) {Test};
    \node[anchor=north west,font=\sffamily\scriptsize,text=ManGrey] at (axis cs:59,175) {Live};

    % traces
    \addplot[AppActual,thick,mark=*,mark size=1pt,domain=0:58,samples=59] {act(x)};
    \addplot[AppBase,very thick,dotted,mark=*,mark size=1pt,mark options={solid},domain=47:65,samples=19] {base(x)};
    \addplot[AppRecA,thick,mark=*,mark size=1pt,domain=47:65,samples=19] {rec(x)};

    % annotations
    \node[man/callout] (a) at (axis cs:20,160) {A};
    \draw[man/arrow,thin] (a) -- (axis cs:22,{act(22)+3});
    \node[man/callout] (b) at (axis cs:44,70) {B};
    \draw[man/arrow,thin] (b) -- (axis cs:50,{base(50)-3});
    \node[man/callout] (c) at (axis cs:54,168) {C};
    \draw[man/arrow,thin] (c) -- (axis cs:57,{rec(57)+3});
    \node[man/callout] (d) at (axis cs:62,82) {D};
    \draw[man/arrow,thin] (d) -- (axis cs:62,100);
    \node[man/callout] (e) at (axis cs:4,140) {E};
    \draw[man/arrow,thin] (e) -- (axis cs:9,{act(9)+2});
  \end{axis}
\end{tikzpicture}

  \caption{Anatomy of the forecast chart (illustrative data)}
  \label{fig:chartAnatomy}
\end{figure}

\begin{table}[H]
\centering
\caption{Lines and regions of the chart (letters as in Figure~\ref{fig:chartAnatomy})}
\label{tab:chartAnatomy}
\small
\begin{tabularx}{\linewidth}{@{}clX@{}}
\toprule
& \textbf{Element} & \textbf{Meaning} \\
\midrule
\Callout{A} & Actual (black, solid, dots) & Units really sold per month; drawn over train and test only \\
\Callout{B} & Base (dotted)              & Forecast of the selected base model; test and live only \\
\Callout{C} & Reconciled (solid, coloured) & One line per selected reconciliation method; test and live only \\
\Callout{D} & Live region                & Forecasts without an actual line -- no real sales to compare with \\
\Callout{E} & Yearly peak                & Seasonal December peak; a good forecast repeats it \\
\bottomrule
\end{tabularx}
\end{table}

Line colours are assigned automatically in a fixed order; the legend \Callout{7} is always the reference for which colour belongs to which method. The legend reads, for example, \UI{Base: ets} and \UI{Reconciled: MinTrace (wls\_struct)}.

\subsection{Chart Toolbar}
\label{sec:uiToolbar}

When the mouse is over the chart, a small toolbar \Callout{10} appears at its top right (Figure~\ref{fig:uiToolbar}). It is provided by the charting library Plotly \cite{Plotly:2026}.

\begin{figure}[H]
  \centering
  \Screenshot[0.55\linewidth]{Images/04UserInterface/ChartToolbar.png}
  \caption{The chart toolbar}
  \label{fig:uiToolbar}
\end{figure}

\begin{table}[H]
\centering
\caption{Chart toolbar buttons}
\label{tab:uiToolbar}
\small
\begin{tabularx}{\linewidth}{@{}lX@{}}
\toprule
\textbf{Button} & \textbf{Function} \\
\midrule
Camera        & Download the current view of the chart as \ac{png} \\
Magnifier     & Zoom mode: drag a rectangle to enlarge it (default mode) \\
Cross arrows  & Pan mode: drag to move the visible range \\
Plus / minus  & Zoom in / out around the centre \\
Autoscale     & Fit all lines into the view \\
House         & Reset the axes to the original view \\
Expand (Streamlit) & Separate button next to the toolbar: show the chart in full screen \\
\bottomrule
\end{tabularx}
\end{table}

\section{Page Menu and Theme}
\label{sec:uiMenu}

The three dots \Callout{11} at the top right of the page open Streamlit's own menu \cite{Streamlit:2026}. For users, two entries matter:
\begin{itemize}
  \item \UI{Settings} -- choose \VAL{Light}, \VAL{Dark} or \VAL{Use system setting}. The chart adjusts its colours automatically (Figure~\ref{fig:uiThemes}).
  \item \UI{Print} -- prints the page including the current chart.
\end{itemize}
On a local installation, a \UI{Deploy} button also appears next to the menu; it is not needed and can be ignored.

\begin{figure}[H]
  \centering
  \Screenshot[0.49\linewidth]{Images/04UserInterface/ChartLight.png}
  \hfill
  \Screenshot[0.49\linewidth]{Images/04UserInterface/ChartDark.png}
  \caption{The same chart in light mode (left) and dark mode (right)}
  \label{fig:uiThemes}
\end{figure}

\section{Messages}
\label{sec:uiMessages}

The dashboard uses Streamlit's standard message boxes. In normal use none of them appears.

\begin{table}[H]
\centering
\caption{Messages the dashboard can show}
\label{tab:uiMessages}
\small
\begin{tabularx}{\linewidth}{@{}llX@{}}
\toprule
\textbf{Type} & \textbf{Beginning of the text} & \textbf{See} \\
\midrule
Red error box    & \texttt{Couldn't find one of the required parquet files} & Section~\ref{sec:tsParquet} \\
Yellow warning   & \texttt{No data found for '\dots'}                        & Section~\ref{sec:tsNoData} \\
Node line suffix & \texttt{No column found for: \dots}                       & Section~\ref{sec:tsNoColumn} \\
Running indicator & (animation, top right)                                   & normal while loading \\
\bottomrule
\end{tabularx}
\end{table}
